\noindent
\subsection{WP6: Community Engagement and Social Impact}

\textbf{Leading Partner:} Eötvös Loránd Tudományegyetem (ELTE BTK)

\textbf{Objectives and Concept:}
WP6 ensures that the research outcomes and digital assets reach a broad audience, fostering strong societal connections. Driven by ELTE's digital humanities expertise, this WP involves local populations, integrates local legends into the XR narratives, and promotes an open science ecosystem where datasets and 3D models become accessible. The ultimate goal is to embed the digital heritage into community identities and measure the resulting social impact.

\textbf{Deliverables / Schedule (from Gantt chart):}
\begin{itemize}
    \item \textbf{D6.1 Community workshops on the relationship between local identity and digital representation:} Year 1 (Month 4-12). These interactive sessions will actively engage residents, local historians, and policymakers to discuss how their regional heritage is being digitized. The objective is to ensure that the 3D assets and virtual narratives accurately and respectfully reflect the community's own cultural self-understanding.
    \item \textbf{D6.2 Co-creation processes: integrating local legends into the XR narrative:} Year 1 (Month 7-12). By utilizing crowdsourced oral histories and regional folklore, the team will enrich the digital reconstructions with authentic, community-driven storytelling. This transforms the VR experience from a sterile architectural walkthrough into a deeply engaging, living historical narrative.
    \item \textbf{D6.3 Preparation of a social impact assessment (SIA):} Year 2 (Month 10-12). This final deliverable will systematically measure and document how the project's digital interventions have influenced local economies, community cohesion, and public attitudes toward heritage conservation. The resulting report will provide empirical data to validate the societal relevance of the "Digital Heritage + XR" approach.
\end{itemize}
